\section{2026 年 11 月核验增补}

本节按论文首次提交时间排列。每项方法附有可核验的外部入选依据。

% METHOD: 2312.02121
\paragraph{gsplat 数学补充：投影、合成与光栅化梯度}
2023.12~\cite{2312_02121}。这篇技术报告把可微高斯光栅化的正向与反向计算逐项展开。正向在高斯中心对透视变换作一阶近似以投影协方差，将二维椭圆按瓦片分组、深度排序并前向合成。反向继续推导合成、投影和三维高斯参数的梯度，便于将公式与轮廓参数和颜色更新对应起来。

它是实现说明，没有独立发布一种新的重建算法或提供性能基准。精确透视投影后的分布也不再是二维高斯，文中的协方差传播采用局部线性近似。作者直接链接的 gsplat 仓库提供光栅化实现、示例与测试，这些公开实现资料是此次收录依据。

\textbf{入选依据} \href{https://github.com/nerfstudio-project/gsplat}{公开实现}已找到与论文对应的代码文件。

% METHOD: 2403.11577
\paragraph{3DGS-Calib：用场景渲染估计多传感器外参与时差}
2024.03~\cite{2403_11577}。3DGS-Calib 将移动平台的激光雷达离散位姿构成连续时间轨迹。给定相对于雷达的候选外参和时间偏移，系统在相应时刻计算相机位姿，并让可微高斯渲染与跨传感器观测共同约束场景和标定量。这样能利用普通场景，不要求专门布置标靶。

作者从 KITTI-360 构建三段序列，分别考察多相机与激光雷达标定。比较中 LCCNet 使用 KITTI 预训练权重，安装几何与 KITTI-360 并不相同。作者还指出当前假设激光雷达点多集中在图像下半部，换成非常规雷达或相机安装时需另行处理。论文作者单位包含华为巴黎与伦敦研究中心。

\textbf{入选依据} 论文署名单位包括 Huawei（\href{https://arxiv.org/abs/2403.11577}{论文原文}）。

% METHOD: 2605.00408
\paragraph{LeGS：用渲染贡献学习高斯的增删与保留}
2026.05~\cite{2605_00408}。LeGS 将自适应密度控制从固定阈值改为可学习的原语决策。训练按单个高斯对全局合成结果的敏感度分配奖励，以保持不变的动作作为回报基准，减少早期不准确的价值估计诱发过度增密。高斯参数本身仍通过常规图像重建损失更新。

原文在 Mip-NeRF 360、Deep Blending 和 Tanks and Temples 报告画质与最终原语数，并与 FastGS 等对照。逐场景结果并非所有指标同步优于对照：例如 Treehill 的 PSNR 小幅降低但感知指标改善。作者仓库包含训练、渲染和评价程序，可追溯到本篇论文。

\textbf{入选依据} \href{https://github.com/aaronnzh/legs}{公开实现}已找到与论文对应的代码文件。

% METHOD: 2605.04509
\paragraph{CoherentRaster：按光场屏幕的子像素布局重新组织高斯渲染}
2026.05~\cite{2605_04509}。光场显示的每个子像素只取某个视点，不必先完整画出全部相机图。CoherentRaster 直接计算参与交织输出的子像素，在邻近视点组复用投影和排序相关属性，再按视点重排 GPU 线程与子像素的对应，缓解交织布局造成的访存分散。

原文用定制 gsplat CUDA 核及统一的 AccuTile 边界和对照比较高分辨率多视图性能。该流程利用相邻视点的一致性，强方向性细节或不连续遮挡会削弱共享近似；实时数字也依赖视点组大小和设备设置。官方仓库发布了训练和渲染包装程序与依赖安装方式。

\textbf{入选依据} \href{https://github.com/sgj0402/coherent-raster}{公开实现}已找到与论文对应的代码文件。仓库名称和实现对应本文；README 里的示例 arXiv 编号与本文不同，以原文编号为准。

% METHOD: 2607.03875
\paragraph{MACRO：用与深度匹配的参考裁剪改进近距离渲染}
2026.07~\cite{2607_03875}。MACRO 将既有 3DGS 的近景图与渲染深度分为多个平面，为不同平面挑选并缩放到相应尺度的参考裁剪，再用随深度变化的注意力掩码修复图像。方法使用既有二维增强模型，无需微调或更改网络架构。

论文在 DL3DV-Closeup 和 MobileClose-10 评价，前者明确隔开远距训练视角与近距测试视角。固定平面数无法适应所有深度复杂度，放大裁剪会丢失高频细节，输入重建的深度错误也会影响处理。作者单位包括 Amazon Prime Video。

\textbf{入选依据} 论文署名单位包括 Amazon（\href{https://arxiv.org/html/2607.03875v2}{论文原文}）。

% METHOD: 2609.03534
\paragraph{TruncGradGS：扩大光栅化反传的有效梯度范围}
2026.09~\cite{2609_03534}。光栅化时，离原语中心较远的像素给出的高斯导数接近零，因此未初始化到表面附近的原语很难移动到正确位置。TruncGradGS 在反向过程中用连续的分段代理替代近零的尾部梯度，使更远的像素也参与更新。前向高斯表示及推理渲染不因这个训练策略改变。

在 Mip-NeRF 360 和新数据集的静态子集，论文分别用随机与 COLMAP 初始化评估 3DGS 和 2DGS，并另评估动态场景。收益在随机初始化下尤其明显。反向加入更多像素与原语配对通常会增加训练时间，因此这种改动改善的是优化可达性，不是无额外训练成本的渲染加速。

\textbf{入选依据} 论文署名单位包括 Dolby Laboratories（\href{https://arxiv.org/abs/2609.03534}{论文原文}）。
